\section{Downstream Outcome Screen}
\label{app:downstream-screen}

This appendix documents the broader downstream screen referenced in Section~\ref{sec:downstream-outcomes}. The screen estimates Borusyak--Jaravel--Spiess imputation event studies for 298 municipality-level outcomes using the strict-BVR non-hybrid sample with never-treated controls, municipality-clustered standard errors, and an event-time window of $-8$ to $+8$. Every reported path is normalized to event time zero by subtracting the event-time-zero estimate from each post-period coefficient. Outcomes span six families: birth and health, education, adult and EJA education, SICONFI log per-capita spending, SICONFI spending shares, and school access and transport.

\textbf{Pretrend screen.} An outcome is labeled \emph{pass} when there are at least three pre-period estimates, no pre-period estimate with $p<0.05$, and a maximum absolute pre-period $t$-statistic no larger than $1.96$. It is labeled \emph{near} when there are at least two pre-period estimates, no more than one pre-period estimate with $p<0.05$, and a maximum absolute pre-period $t$-statistic no larger than $2.58$. All other outcomes are labeled \emph{fail}. Of the 298 estimated outcomes, 42 pass the screen, 20 are near, and 236 fail. The main text plots only pass-pretrend outcomes in the political-economy mechanism subsections, and reports the near-pretrend teen-mother outcome with an explicit measurement-visibility caveat.

\textbf{Ranking and run summary.} The master ranking table, the stacked event-study estimates, and the augmented diagnostics file are stored under \texttt{resources/tables/downstream\_bjs\_ref0\_nonhybrid/}. The corresponding estimator outputs are stored under \texttt{resources/did/downstream\_bjs\_ref0\_nonhybrid/}. Combined figures for each outcome family are stored under \texttt{resources/figures/downstream\_bjs\_ref0\_nonhybrid/}.

\textbf{Multiple-testing discipline.} Because the screen estimates a large number of outcomes, the pass-pretrend set should not be interpreted as a battery of independent confirmatory tests. Conditioning on a pre-period diagnostic can also bias post-period contrasts when the pre-period test is noisy. I therefore treat the pass-pretrend outcomes as mechanism candidates rather than as final causal estimates, and the main-text mechanism story relies on the coherence of the surviving pattern across spending, school flow, and selected birth and immunization outcomes rather than on any single specification.
